\chapter{Default EIRENE tallies, and selected Modules}\label{sec.tables-modules}
\section{Tables of EIRENE tallies}\label{sec.tables}

The following three tables comprise the EIRENE pre-programmed default ``tallies".
(Further tallies can be added via the problem specific interfacing or user routines).

The term ``tally" is adopted from neutron transport applications, a
more precise terminology would be ``response", see
\cref{sec3.2}.  By abuse of language we also refer to input data, such as plasma profiles
defined on the computational grid, as tallies, because in post-processing these can printed, plotted
and be algebraically mixed with the output tallies (responses) to form more general output quantities.
The three tables given below list the (volumetric) input tallies, (volumetric, cell averaged) output tallies,
and the surface averaged output tallies, respectively.

Because of the new type of particle introduced
during 2002 (photon gas, ITYP=0) the numbering of tallies has
changed significantly. Therefore the tables are given below once
for the current EIRENE version (2002 and younger), and once for
the versions older than 2001. A further change in tally numbering and naming conventions
has been carried out in early 2014, and this affects the surface sputter tallies, (and those
surface tallies in the list after the sputter tallies.
These latter revisions are described in \cref{tab3a}, and only for that section of the surface tally list
that has been affected.

\

Each \textbf{volume averaged response} is a spatial function,
averaged over a grid cell, hence: a piecewise constant function, which is constant in each grid
cell ICELL:
\\
ICELL=1,NRAD.  %
\\
Default volumetric particle, energy and momentum source tallies are
split by the \textbf{type} (but not the individual species index of
the incident test particle) involved in a collision process: photons
(Phs.), atoms (Ats.), molecules (Mls.), test ions (T.I.). This
permits scaling of these sources with the factors FPHOT, FATM, FMOL,
FION, to eliminate statistical errors from the particle balances
(see NLSCL option, input block 1). They are further split by both:
the type (also: electrons (El.) and the species index of the ``receiving''
component, for which the source/sink term would arise in a velocity
space averaged balance equation (with the sign as specified).
Any further tallies (e.g. by incident species, or particular collision process resolved,
or energy (spectrally) resolved, etc.)
have to be implemented in the 3rd party (user) or the problem specific
interfacing routines.

Tally naming convention:
\begin{list}{ }{}
\item[{Particle source rates}]

name: P[A][BC], with\\
A: type of incident particle, A = A (atoms), M (molecules), I (test ions),  PH (photons), P (bulk particle)\\
BC: type of produced particle, BC = AT (atoms), ML (molecules) , IO (test ions), PHT (photons), EL (electrons) , PL (bulk particles)\\
The tallies are resolved by species index of produced particle, and are to be scaled (NLSCL option) by type of incident test particle.
\item[{Energy source rates}]

name: E[A][BC], analogue to particle source rates\\
but not resolved wrt. species index of secondary particles.
\item[{Momentum source rates}]

name: M[A][BC], analogue to particle source rates\\
only available for [A] = A, M, I, PH  and for [BC]=PL,
i.e. currently default tallies are only available for momentum sources/sinks for background (plasma) particles caused by
test particles, resolved by type of test particle (name of tally) and by species of bulk particle (1st index in array).
\end{list}

\textbf{Surface averaged responses} are also spatial
functions, averaged over surface segments, hence:
piecewise constant on each surface ISURF or surface segment:
\\
ISURF=1,NLIMI.
\\
and with some spatial resolution possible on standard mesh surfaces
\\
ISURF=NLIM,NLIM+NSTSI
\\
Emitted particle and energy surface averaged fluxes are split by
the typ of the incident particle: photons (Phs.), atoms (Ats.),
molecules (Mls.), test ions (T.I.) or bulk ions (B.I.). This
permits scaling of these fluxes with the factors FPHOT, FATM,
FMOL, FION, to eliminate statistical errors from the particle
balances (see NLSCL option, input block 1).

\medskip
Naming conventions for emitted surface fluxes are analogue to those for volumetric source rates,
i.e. PRF[A][BC],  SPT[A][BC] with A and BC having the same meaning as described above for volume tallies.\\
PRF... are fluxes reflected or emitted due to incident particle fluxes, of type [A], SPT[A][BC] are sputtered fluxes (additional
fluxes of surface material).

\textbf{Note:} The naming convention of \textbf{surface sputter tallies SPT...} differs in older versions of EIRENE.
The full resolution wrt. to incident type and emitted species is only implemented in EIRENE revisions svn 360 and later (Jan. 2014), see \cref{tab3a}.

In all older versions only much fewer sputter tallies SPT[BC] were available, i.e. only one type-identifier in the name of these ``sputter tallies".

And [BC] (and the corresponding species index) could either mean:\\
i)  resolution wrt. incident flux species (summing over all sputtered flux species and types) (same as in incident surface flux tallies), or \\
ii) resolution wrt. sputtered flux (summing over all incident species and types), \\
depending on code version. Type ii) convention was used in EIRENE versions operated at ITER-IO (e.g. in SOLPS4.3), all other
EIRENE versions had convention i).

The extension of Jan. 2014 now combines i) and ii) by providing the resolution wrt. emitted type and species (as ii) did) and
simultaneously keeping the resolution wrt. incident type (i), as needed for proper scaling of particle balances (see option
NLSCL, input block 1).


\medskip
If temporal
resolution has been requested (input block 13), then NSTSI is increased
by one: NSTSIP=NSTSI+1, and the last tally I2=NLIM+NSTSIP corresponds
to the ``Time-Surface" (the census array) then.

%\newpage

\subsection{Current status, incl. photon gas tallies (Eirene\texorpdfstring{$_{2002}$}{2002} and
younger)}\label{sec.tables-new}
\begin{table}[h]
\caption{\textbf{Input Tallies for Background,
Input, Module: COMUSR}}
\bigskip
\label{tab1} {\footnotesize
\begin{tabular} {|r|l|l|l|l|c|} \hline

 \textbf{No} & \textbf{Name} & \textbf{Macroscopic quantity}
 &  \textbf{1.Dim.}
  & \textbf{Units} & \textbf{Estim.}   \\ \hline

 -1 & TEIN & Plasma Temperature, Electrons & 1 & \si{\electronvolt} & /  \\ \hline
 -2 & TIIN & Plasma Temperature, Bulk Ions & NPLSTI & \si{\electronvolt} & /  \\ \hline
 -3 & DEIN$^*$ & Plasma Density, Electrons & 1 & \si{\per\cubic\centi\metre} & /  \\ \hline
 -4 & DIIN & Plasma Density, Bulk Ions & NPLS & \si{\per\cubic\centi\metre} & / \\ \hline
 -5 & VXIN & Plasma Drift Velocity, x-component, Bulk Ions&NPLSV
     & \si{\centi\metre\per\second} & /                                           \\ \hline
 -6 & VYIN & Plasma Drift Velocity, y-component, Bulk Ions&NPLSV
     & \si{\centi\metre\per\second} & /                                           \\ \hline
 -7 & VZIN & Plasma Drift Velocity, z-component, Bulk Ions&NPLSV
     & \si{\centi\metre\per\second} & /                                           \\ \hline
 -8 & BXIN & Magnetic field unit vector, x-component & 1 & / & /
    \\ \hline
 -9 & BYIN & Magnetic field unit vector, y-component & 1 & / & /
    \\ \hline
 -10 & BZIN & Magnetic field unit vector, z-component & 1 & / & /
    \\ \hline
 -11 & BFIN & Magnetic field strength & 1 &  Tesla & /
    \\ \hline
 -12 & ADIN & Additional input tally & NAIN & see INFCOP & /
    \\ \hline
 -13 & EDRIFT$^*$ & Kinetic energy in drift motion, Bulk Ions &NPLS & \si{\electronvolt} & /
    \\ \hline
 -14 & VOL & Zone Volume & 1 & \si{\cubic\centi\metre} & /
    \\ \hline
 -15 & WGHT $^{**}$ & unused, importance fct.  & NSPCMC & $ 1 $ & /
 \\ \hline
-16  &   BXPERP $^*$  &  unused   &       1 &     &    /
   \\ \hline
-17  &   BYPERP $^*$  &  unused  &      1   &     &   /
    \\ \hline
-18   &      EXIN   &       electric field unit vector, x component &  1  &   &  /
    \\ \hline
-19   &      EYIN   &      electric field unit vector, y component&  1  &   &  /
     \\ \hline
-20   &      EZIN   &    electric field unit vector, z component  &  1  &   &  /
     \\ \hline
-21   &      EFIN   &      electric field strength  & 1  &         V/CM    &/
     \\ \hline
-22   &      POT   &      electric potential  & 1  &         V    &/
     \\ \hline\hline
d1    &      BVIN $^*$ &      flow velocity parallel to B field  & NPLSV  &    CM/S    &/
     \\ \hline
d2    &      PARMOM $^*$ &    parallel to B flow momentum   & NPLS  &  G*CM/S    &/
     \\ \hline
d3    &      TEDTEDX $^{**}$ &    unused  & 1  &  --    &/
     \\ \hline
d4    &      TEDTEDY $^{**}$ &    unused  & 1  &  --    &/
     \\ \hline
d5    &      TEDTEDZ $^{**}$ &    unused  & 1  &  --    &/
     \\ \hline
\end{tabular}%
}
\end{table}

\medskip

\textbf{Note:}
We distinguish between so called primary input tallies,
and derived input tallies, the latter being  indicated by $^*$.
The latter are automatically computed from the former by EIRENE.


Primary input tallies are set or transferred from external data sources via subr. PLASMA.f,
derived input tallies are set via subroutine PLASMA$\_$DERIV.f.
Storage for both kinds of tallies is allocated in module COMUSR.f

\begin{itemize}

\item NPLSTI:  number of different bulk particle (bulk ion) temperatures in a run.
See input block 5. Default: NPLSTI=NPLS

\item NPLSV:   number of different bulk particle (bulk ion) flow velocities in a run.
See input block 5. Default: NPLSV=NPLS

\item DEIN and EDRIFT are ``derived" input tallies,
internally computed from the ion densities (quasi-neutrality), and
the flow field, respectively.
\item Tallies d1 -- d5 are further ``internal" derived quantities, currently not selectable for printout
nor for graphics.
\item  B-field and E-field tallies contain redundancy. Storage can be reduced
\item  WGHT space species  weight window arrays.  Currently not in use

\item  Magnetic field input is needed only for photon transport with Zeeman splitted line shapes, or
in case of charged particle transport (Subr. FOLION). Only quite rudimentary preprogrammed input options
for magnetic field are available, usually it is expected to receive B-field information from external files,
see INDPRO(5) options for magnetic field tallies in input block 5 (\cref{sec2.5}).
\item  For electric fields (needed only for charged particle transport) only input via external files
is possible. Default: $\vec{E} = 0$
\end{itemize}



\newpage

\clearpage
{\footnotesize
%\begin{longtable}[t]{|r|l|l|l|l|c|}
\begin{longtable}{|r|l|l|l|l|c|}
\caption{\textbf{Volume Averaged Tallies, Output,
Module: CESTIM}}
\label{tab2}
\\
\hline
\textbf{No} & \textbf{Name} & \textbf{Macroscopic quantity}
& \textbf{1.Dim.}
& \textbf{Units} & \textbf{Estim.} \\
\hline\hline
\endfirsthead
\caption[]{(continued)}\\
\hline
 \textbf{No} & \textbf{Name} & \textbf{Macroscopic quantity}
 & \textbf{1.Dim.}
 & \textbf{Units} & \textbf{Estim.} \\
\hline\hline
\endhead
 1 & PDENA & Particle density, Atoms & NATM & \si{\per\cubic\centi\metre} & TD
 \\ \hline
 2 & PDENM & Particle density, Molecules & NMOL & \si{\per\cubic\centi\metre} & TD
 \\ \hline
 3 & PDENI & Particle density, Test Ions & NION  & \si{\per\cubic\centi\metre} & TD
 \\ \hline
 4 & PDENPH & Particle density, Photons & NPHOT  & \si{\per\cubic\centi\metre} & TD
 \\ \hline
 5 & EDENA & Energy density, Atoms & NATM       & \si{\electronvolt}*\si{\per\cubic\centi\metre} & TE
 \\ \hline
 6 & EDENM & Energy density, Molecules & NMOL   & \si{\electronvolt}*\si{\per\cubic\centi\metre} & TE
 \\ \hline
 7 & EDENI & Energy density, Test Ions & NION   & \si{\electronvolt}*\si{\per\cubic\centi\metre} & TE
 \\ \hline
 8 & EDENPH & Energy density, Photons & NPHOT   & \si{\electronvolt}*\si{\per\cubic\centi\metre} & TE
 \\ \hline
 9 & PAEL & Particle Source (Electrons) from atom-plasma coll.
   & 1 & \si{\ampere}* \si{\per\cubic\centi\metre} & TPS  \\ \hline
 10 & PAAT & Particle Source (Atoms) from atom-plasma coll.
   & NATM & \si{\ampere}*\si{\per\cubic\centi\metre} & TPS \\ \hline
 11 & PAML & Particle Source (Molecules) from atom-plasma coll.
   & NMOL & \si{\ampere}*\si{\per\cubic\centi\metre} & TPS \\  \hline
 12 & PAIO & Particle Source (Test Ions) from atom-plasma coll.
    & NION & \si{\ampere}*\si{\per\cubic\centi\metre} & TPS  \\  \hline
 13 & PAPHT & Particle Source (Photons) from atom-plasma coll.
    & NPHOT & \si{\ampere}*\si{\per\cubic\centi\metre} & TPS  \\  \hline
 14 & PAPL & Particle Source (Bulk Ions) from atom-plasma coll.
    & NPLS & \si{\ampere}*\si{\per\cubic\centi\metre} & TPS  \\  \hline
 15 & PMEL &Particle Source (Electrons) from molecule-plasma coll.
    & 1 & \si{\ampere}*\si{\per\cubic\centi\metre} & TPS     \\  \hline
 16 & PMAT &Particle Source (Atoms) from molecule-plasma coll.
    & NATM &\si{\ampere}*\si{\per\cubic\centi\metre} & TPS  \\  \hline
 17 & PMML &Particle Source (Molecules) from molecule-plasma coll.
    & NMOL &\si{\ampere}*\si{\per\cubic\centi\metre} & TPS   \\  \hline
 18 & PMIO &Particle Source (Test Ions) from molecule-plasma coll.
    & NION &\si{\ampere}*\si{\per\cubic\centi\metre} & TPS     \\   \hline
 19 & PMPHT &Particle Source (Photons) from molecule-plasma coll.
    & NPHOT &\si{\ampere}*\si{\per\cubic\centi\metre} & TPS     \\   \hline
 20 & PMPL &Particle Source (Bulk Ions) from molecule-plasma coll.
    & NPLS &\si{\ampere}*\si{\per\cubic\centi\metre} & TPS    \\    \hline
 21 & PIEL &Particle Source (Electrons) from test ion-plasma coll.
    & 1 &\si{\ampere}*\si{\per\cubic\centi\metre} & TPS       \\   \hline
 22 & PIAT &Particle Source (Atoms) from test ion-plasma coll.
    & NATM &\si{\ampere}*\si{\per\cubic\centi\metre} & TPS    \\  \hline
 23 & PIML &Particle Source (Molecules) from test ion-plasma coll.
    & NMOL &\si{\ampere}*\si{\per\cubic\centi\metre} & TPS    \\  \hline
 24 & PIIO &Particle Source (Test Ions) from test ion-plasma coll.
    & NION &\si{\ampere}*\si{\per\cubic\centi\metre} & TPS    \\   \hline
 25 & PIPHT &Particle Source (Photons) from test ion-plasma coll.
    & NPHOT &\si{\ampere}*\si{\per\cubic\centi\metre} & TPS    \\   \hline
 26 & PIPL &Particle Source (Bulk Ions) from test ion-plasma coll.
    & NPLS &\si{\ampere}*\si{\per\cubic\centi\metre} & TPS    \\   \hline
 27 & PPHEL &Particle Source (Electrons) from photon-plasma coll.
    & 1 &\si{\ampere}*\si{\per\cubic\centi\metre} & TPS       \\   \hline
 28 & PPHAT &Particle Source (Atoms) from photon-plasma coll.
    & NATM &\si{\ampere}*\si{\per\cubic\centi\metre} & TPS    \\  \hline
 29 & PPHML &Particle Source (Molecules) from photon-plasma coll.
    & NMOL &\si{\ampere}*\si{\per\cubic\centi\metre} & TPS    \\  \hline
 30 & PPHIO &Particle Source (Test Ions) from photon-plasma coll.
    & NION &\si{\ampere}*\si{\per\cubic\centi\metre} & TPS    \\   \hline
 31 & PPHPHT &Particle Source (Photons) from photon-plasma coll.
    & NPHOT &\si{\ampere}*\si{\per\cubic\centi\metre} & TPS    \\   \hline
 32 & PPHPL &Particle Source (Bulk Ions) from photon-plasma coll.
    & NPLS &\si{\ampere}*\si{\per\cubic\centi\metre} & TPS    \\
\hline
\pagebreak
 33 & EAEL &Energy Source (Electrons) from atom-plasma coll.
    & 1 & \si{\watt}*\si{\per\cubic\centi\metre} & T22      \\  \hline
 34 & EAAT &Energy Source (Atoms) from atom-plasma coll.
    & 1 & \si{\watt}*\si{\per\cubic\centi\metre} & T23      \\   \hline
 35 & EAML &Energy Source (Molecules) from atom-plasma coll.
    & 1 & \si{\watt}*\si{\per\cubic\centi\metre} & T24      \\    \hline
 36 & EAIO &Energy Source (Test Ions) from atom-plasma coll.
    & 1 & \si{\watt}*\si{\per\cubic\centi\metre} & T25      \\ \hline
 37 & EAPHT &Energy Source (Photons) from atom-plasma coll.
    & 1 & \si{\watt}*\si{\per\cubic\centi\metre} & T25      \\ \hline
 38 & EAPL &Energy Source (Bulk Ions) from atom-plasma coll.
    & 1 & \si{\watt}*\si{\per\cubic\centi\metre} & T26   \\  \hline
 39 & EMEL &Energy Source (Electrons) from molecule-plasma coll.
    & 1 & \si{\watt}*\si{\per\cubic\centi\metre} & T27   \\   \hline
 40 & EMAT &Energy Source (Atoms) from molecule-plasma coll.
    & 1 & \si{\watt}*\si{\per\cubic\centi\metre} & T28  \\   \hline
 41 & EMML &Energy Source (Molecules) from molecule-plasma coll.
    & 1 & \si{\watt}*\si{\per\cubic\centi\metre} & T29  \\   \hline
 42 & EMIO &Energy Source (Test Ions) from molecule-plasma coll.
    & 1 & \si{\watt}*\si{\per\cubic\centi\metre} & T30  \\    \hline
 43 & EMPHT &Energy Source (Photons) from molecule-plasma coll.
    & 1 & \si{\watt}*\si{\per\cubic\centi\metre} & T30  \\    \hline
 44 & EMPL & Energy Source (Bulk Ions) from molecule-plasma coll.
    & 1& \si{\watt}*\si{\per\cubic\centi\metre} & T31  \\  \hline
 45 & EIEL & Energy Source (Electrons) from test ion-plasma coll.
    & 1& \si{\watt}*\si{\per\cubic\centi\metre} & T32   \\  \hline
 46 & EIAT & Energy Source (Atoms) from test ion-plasma coll.
    & 1& \si{\watt}*\si{\per\cubic\centi\metre} & T33   \\  \hline
 47 & EIML & Energy Source (Molecules) from test ion-plasma coll.
    & 1& \si{\watt}*\si{\per\cubic\centi\metre} & T34   \\ \hline
  48 & EIIO & Energy Source (Test Ions) from test ion-plasma coll.
    & 1& \si{\watt}*\si{\per\cubic\centi\metre} & T35   \\  \hline
 49 & EIPHT & Energy Source (Photons) from test ion-plasma coll.
    & 1& \si{\watt}*\si{\per\cubic\centi\metre} & T35   \\  \hline
 50 & EIPL & Energy Source (Bulk Ions) from test ion-plasma coll.
    & 1& \si{\watt}*\si{\per\cubic\centi\metre} & T36   \\  \hline
 51 & EPHEL & Energy Source (Electrons) from photon-plasma coll.
    & 1& \si{\watt}*\si{\per\cubic\centi\metre} & T33   \\  \hline
 52 & EPHAT & Energy Source (Atoms) from photon-plasma coll.
    & 1& \si{\watt}*\si{\per\cubic\centi\metre} & T33   \\  \hline
 53 & EPHML & Energy Source (Molecules) from photon-plasma coll.
    & 1& \si{\watt}*\si{\per\cubic\centi\metre} & T34   \\  \hline
 54 & EPHIO & Energy Source (Test Ions) from photon-plasma coll.
    & 1& \si{\watt}*\si{\per\cubic\centi\metre} & T35   \\  \hline
 55 & EPHPHT & Energy Source (Photons) from photon-plasma coll.
    & 1& \si{\watt}*\si{\per\cubic\centi\metre} & T35   \\  \hline
 56 & EPHPL & Energy Source (Bulk Ions) from photon-plasma coll.
    & 1& \si{\watt}*\si{\per\cubic\centi\metre} & T36   \\ \hline
 57 & ADDV & Additional volume av. Tally, Track-length estimated
    & NADV & see UPTUSR & TRL      \\  \hline
 58 & COLV & Additional volume av. Tally, Collision estimated
    & NCLV & see UPCUSR & COL        \\  \hline
 59 & SNPV & Additional volume av. Tally, Snapshot estimated
    & NSNV & see UPSUSR & SNP         \\  \hline
 60 & COPV & Tallies for coupling to ext. code, see: Subr. UPTCOP
    & NCPV & see UPTCOP  & COP  \\  \hline
 61 & BGKV &Volume averaged tallies for \acs{BGK}-self-collision terms
    & NBGV & see BGK & T41        \\  \hline
 62 & ALGV &Algebraic expression in volume averaged tallies
    & NALV & Input & ALG        \\  \hline
 63 & PGENA & Generation limit, Atoms & NATM & \si{\ampere}*\si{\per\cubic\centi\metre} & CG
 \\ \hline
 64 & PGENM & Generation limit, Molecules & NMOL & \si{\ampere}*\si{\per\cubic\centi\metre} & CG
 \\
\hline
\pagebreak
 65 & PGENI & Generation limit, Test ions & NION & \si{\ampere}*\si{\per\cubic\centi\metre} & CG
 \\ \hline
 66 & PGENPH & Generation limit, Photons & NPHOT & \si{\ampere}*\si{\per\cubic\centi\metre} & CG
 \\ \hline
 67 & EGENA & dito, Energy flux, Atoms & NATM & \si{\watt}*\si{\per\cubic\centi\metre} & CG
 \\ \hline
 68 & EGENM & dito, Energy flux, Molecules & NMOL & \si{\watt}*\si{\per\cubic\centi\metre} & CG
 \\ \hline
 69 & EGENI & dito, Energy flux, Test ions & NION & \si{\watt}*\si{\per\cubic\centi\metre} & CG
 \\ \hline
 70 & EGENPH & dito, Energy flux, Photons & NPHOT & \si{\watt}*\si{\per\cubic\centi\metre} & CG
 \\ \hline
 71 & VGENA & dito, momentum flux, Atoms & NATM & as MAPL & CG
 \\ \hline
 72 & VGENM & dito, momentum flux, Molecules & NMOL & as MMPL & CG
 \\ \hline
 73 & VGENI & dito, momentum flux, Test ions & NION & as MIPL & CG
  \\ \hline
 74 & VGENPH & dito, momentum flux, Photons & NPHOT & as MPHPL & CG
\\ \hline
75 & PPAT & primary particle sources rate (Atoms) from plasma & NATM &
\si{\ampere}*\si{\per\cubic\centi\metre} & C1
\\ \hline
76 & PPML & primary particle sources rate (Molecules) & NMOL &
\si{\ampere}*\si{\per\cubic\centi\metre} & C1
\\ \hline
77 & PPIO & primary particle sources rate (Test ions) & NION &
\si{\ampere}*\si{\per\cubic\centi\metre} & C1
\\ \hline
78 & PPPHT & primary particle sources rate (Photons) & NPHOT &
\si{\ampere}*\si{\per\cubic\centi\metre} & C1
\\ \hline
79 & PPPL & primary particle sources rate (Bulk Ions) & NPLS &
\si{\ampere}*\si{\per\cubic\centi\metre} & C1
\\ \hline
80 & EPAT & primary energy sources rate, Atoms & 1 & \si{\watt}*\si{\per\cubic\centi\metre}
& C1
\\ \hline
81 & EPML & primary energy sources rate, Molecules & 1 &
\si{\watt}*\si{\per\cubic\centi\metre} & C1
\\ \hline
82 & EPIO & primary energy sources rate, Test ions & 1 &
\si{\watt}*\si{\per\cubic\centi\metre} & C1
\\ \hline
83 & EPPHT & primary energy sources rate, Photons & 1 &
\si{\watt}*\si{\per\cubic\centi\metre} & C1
\\ \hline
84 & EPPL & primary energy sources rate, Bulk Ions & 1 &
\si{\watt}*\si{\per\cubic\centi\metre} & C1
\\ \hline
85 & VXDENA & momentum density, x-direction, Atoms & NATM &
\si{\gram\centi\metre\per\second\per\cubic\centi\metre} & T1
\\ \hline
86 & VXDENM & momentum density, x-direction, Molecules & NMOL &
\si{\gram\centi\metre\per\second\per\cubic\centi\metre} & T1
\\ \hline
87 & VXDENI & momentum density, x-direction, Test Ions & NION &
\si{\gram\centi\metre\per\second\per\cubic\centi\metre} & T1
\\ \hline
88 & VXDENPH & momentum density, x-direction, Photons & NPHOT &
\si{\gram\centi\metre\per\second\per\cubic\centi\metre} & T1
\\ \hline
89 & VYDENA & momentum density, y-direction, Atoms & NATM &
\si{\gram\centi\metre\per\second\per\cubic\centi\metre} & T1
\\ \hline
90 & VYDENM & momentum density, y-direction, Molecules & NMOL &
\si{\gram\centi\metre\per\second\per\cubic\centi\metre} & T1
\\ \hline
91 & VYDENI & momentum density, y-direction, Test Ions & NION &
\si{\gram\centi\metre\per\second\per\cubic\centi\metre} & T1
\\ \hline
92 & VYDENPL & momentum density, y-direction, Photons & NPHOT &
\si{\gram\centi\metre\per\second\per\cubic\centi\metre} & T1
\\ \hline
93 & VZDENA & momentum density, z-direction, Atoms & NATM &
\si{\gram\centi\metre\per\second\per\cubic\centi\metre} & T1
\\ \hline
94 & VZDENM & momentum density, z-direction, Molecules & NMOL &
\si{\gram\centi\metre\per\second\per\cubic\centi\metre} & T1
\\ \hline
95 & VZDENI & momentum density, z-direction, Test Ions & NION &
\si{\gram\centi\metre\per\second\per\cubic\centi\metre} & T1
\\ \hline
96 & VZDENPL & momentum density, z-direction, Photons & NPHOT &
\si{\gram\centi\metre\per\second\per\cubic\centi\metre} & T1 \\
\hline \pagebreak 97 & MAPL & parallel momentum source rate (Bulk Ions) &
NPLS & \si{\ampere}*\si{\gram\centi\metre\per\second\per\cubic\centi\metre} & TPM
\\ \hline
98 & MMPL & parallel momentum source rate, (Bulk Ions) & NPLS &
\si{\ampere}*\si{\gram\centi\metre\per\second\per\cubic\centi\metre} & TPM
\\ \hline
99 & MIPL & parallel momentum source rate, (Bulk Ions) & NPLS &
\si{\ampere}*\si{\gram\centi\metre\per\second\per\cubic\centi\metre} & TPM
\\ \hline
100 & MPHPL & parallel momentum sources rate, (Bulk Ions) & NPLS &
\si{\ampere}*\si{\gram\centi\metre\per\second\per\cubic\centi\metre} & TPM
\\ \hline
\end{longtable}
}
\textbf{Note:}
$\underline{Momentum  \ density}$  is the directional momentum weighted velocity space integral, i.e. it is a (directional) particle
flux (density), for the x,y and z directions, respectively, and then multiplied by the particle mass. (Relation to particle current??,
check with neutron transport terminology)

$\underline{Source} $ means: if the sign is positive, it is a gain
for the specified type of particles; if it is negative, it is a loss.
The specified "type of particle" is coded by the last two letters of the name of a tally, AT, ML, IO, EL, PL, PH, respectively,
as well as by the corresponding ``1st dimension" (species)  NATM, NMOL, NION, NPHOT, NPLS and ``1" (one) for electrons..
e.g. MAPL:  parallel momentum source/sink for bulk particle (plasma background), due to atom plasma interactions. The parallel direction
is determined by the magnetic field unit vector (input tallies BXIN,...).

\bigskip
\textbf{Estimators} EIRENE resorts, by default, to tracklength
estimators \cref{eq1.11}. Default tallies are updated (``scoring")
in routine UPDATE. All default estimators are constant within a cell
$m$, i.e. depend only upon the cell index $m$ but not on the
position $\vr$ in the cell: $g_t(s) = const(m)$. Hence they also do
not depend on the position $s$ along the track within a cell. UPDATE
calls templates UPTUSR, in which any further quantity can be scored
by programming the path integral of any function $g_t(s)$, see
\cref{sec3.2}.

In some instances still collision estimators \cref{eq1.9} are
employed but we are gradually removing them (and plan keep them in
the code only to provide independent checks for code-verification
runs.)
\begin{itemize}
\item[TD]  Tracklength, particle density. $g_t= 1/v$, $v=VEL$ the
test particle's velocity. Note that the path-integral of $g_t$ along
a trajectory in a cell is equal to the time spend by that history in
a cell.
\item[TE]  Tracklength, (kinetic) energy density. In case of photon tallies, the energy used is $E0 = h\nu$
\item[TPS] Tracklength, particle source rate
\item[TPM] Tracklength, parallel (to B-field) momentum source/sink. The (vectorial) momentum exchange rate
$\vec{\Delta Mom}$ is projected to the B-field (input volume tallies Bxin, Byin, Bzin). I.e. if $e_{\vec B}$ is the unit
vector given by input tallies Bxin, Byin, Bzin,
then the scalar product $\langle e_{\vec B},
\vec{\Delta Mom}\rangle$ is used for scoring.
\item[CG]  Collision estimator for generation limit (for fluid limit) tallies. These tallies contain particle, momentum and energy fluxes
(per cell or per \si{\cubic\centi\metre}), when particles are removed from the fully kinetic tracing, due to e.g. transition into a
diffusive or other continuum regime. These fluxes are scored in COLLIDE.F, i.e. by collision estimators by nature.
\end{itemize}


%\newpage
%\begin{table}[t]
%\caption \textbf{Surface Averaged Tallies,
%Output, Module: CESTIM}
%\bigskip
%\label{tab3}

\clearpage
{\footnotesize
\begin{longtable}{|r|l|l|l|l|c|}
\caption{\textbf{Surface Averaged Tallies,
Output, Module: CESTIM}}
\label{tab3}
\\
\hline

 \textbf{No} & \textbf{Name} & \textbf{Macroscopic quantity}
&
   \textbf{1.Dim.}
  & \textbf{Units} & \textbf{Estim.}   \\
\hline\hline
\endfirsthead
\caption[]{(continued)}\\
\hline
 \textbf{No} & \textbf{Name} & \textbf{Macroscopic quantity}
 & \textbf{1.Dim.}
 & \textbf{Units} & \textbf{Estim.} \\
\hline\hline
\endhead
 1 & POTAT &Particle Flux, incident, Atoms&NATM
    & \si{\ampere} & T/C  \\ \hline
 2 & PRFAAT &Particle Flux, emitted, Ats.$\Rightarrow$Atoms&NATM
    & \si{\ampere} & T/C  \\  \hline
 3 & PRFMAT &Particle Flux, emitted, Mls.$\Rightarrow$Atoms&NATM
    & \si{\ampere} & T/C  \\  \hline
 4 & PRFIAT &Particle Flux, emitted, T.I.$\Rightarrow$Atoms&NATM
    & \si{\ampere} & T/C  \\  \hline
 5 & PRFPHAT &Particle Flux, emitted, Pht.$\Rightarrow$Atoms&NATM
    & \si{\ampere} & T/C  \\  \hline
 6 & PRFPAT &Particle Flux, emitted, B.I.$\Rightarrow$Atoms&NATM
    & \si{\ampere} & T/C  \\  \hline
 7 & POTML &Particle Flux, incident, Molecules&NMOL
    & \si{\ampere} & T/C  \\  \hline
 8 & PRFAML &Particle Flux, emitted, Ats.$\Rightarrow$Molecules&NMOL
    & \si{\ampere} & T/C  \\   \hline
 9 & PRFMML &Particle Flux, emitted, Mls.$\Rightarrow$Molecules&NMOL
    & \si{\ampere} & T/C  \\   \hline
 10 & PRFIML &Particle Flux, emitted, T.I.$\Rightarrow$Molecules&NMOL
    & \si{\ampere} & T/C  \\   \hline
 11 & PRFPHML &Particle Flux, emitted, Pht.$\Rightarrow$Molecules&NMOL
    & \si{\ampere} & T/C  \\   \hline
 12 & PRFPML &Particle Flux, emitted, B.I.$\Rightarrow$Molecules&NMOL
    & \si{\ampere} & T/C  \\   \hline
 13 & POTIO &Particle Flux, incident, Test Ions&NION
    & \si{\ampere} & T/C  \\    \hline
 14 & PRFAIO &Particle Flux, emitted, Ats.$\Rightarrow$Test Ions&NION
    & \si{\ampere} & T/C  \\    \hline
 15 & PRFMIO &Particle Flux, emitted, Mls.$\Rightarrow$Test Ions&NION
    & \si{\ampere} & T/C  \\    \hline
 16 & PRFIIO &Particle Flux, emitted, T.I.$\Rightarrow$Test Ions&NION
    & \si{\ampere} & T/C  \\    \hline
 17 & PRFPHIO &Particle Flux, emitted, Pht.$\Rightarrow$Test Ions&NION
    & \si{\ampere} & T/C  \\    \hline
 18 & PRFPIO &Particle Flux, emitted, B.I.$\Rightarrow$Test Ions&NION
    & \si{\ampere} & T/C  \\    \hline
 19 & POTPHT &Particle Flux, incident, Photons&NPHOT
    & \si{\ampere} & T/C  \\    \hline
 20 & PRFAPHT &Particle Flux, emitted, Ats.$\Rightarrow$Photons&NPHOT
    & \si{\ampere} & T/C  \\    \hline
 21 & PRFMPHT &Particle Flux, emitted, Mls.$\Rightarrow$Photons&NPHOT
    & \si{\ampere} & T/C  \\    \hline
 22 & PRFIPHT &Particle Flux, emitted, T.I.$\Rightarrow$Photons&NPHOT
    & \si{\ampere} & T/C  \\    \hline
 23 & PRFPHPHT &Particle Flux, emitted, Pht.$\Rightarrow$Photons&NPHOT
    & \si{\ampere} & T/C  \\    \hline
 24 & PRFPPHT &Particle Flux, emitted, B.I.$\Rightarrow$Photons&NPHOT
    & \si{\ampere} & T/C  \\    \hline
 25 & POTPL &Particle Flux, incident, Bulk Ions&NPLS
    & \si{\ampere} & T/C  \\
\hline
\pagebreak
 26 & EOTAT &Energy Flux, incident, Atoms&NATM
    & \si{\watt} & T/C  \\   \hline
 27 & ERFAAT &Energy Flux, emitted, Ats.$\Rightarrow$Atoms&NATM
    & \si{\watt} & T/C  \\   \hline
 28 & ERFMAT &Energy Flux, emitted, Mls.$\Rightarrow$Atoms&NATM
    & \si{\watt} & T/C  \\   \hline
 29 & ERFIAT &Energy Flux, emitted, T.I.$\Rightarrow$Atoms&NATM
    & \si{\watt} & T/C  \\   \hline
 30 & ERFPHAT &Energy Flux, emitted, Pht.$\Rightarrow$Atoms&NATM
    & \si{\watt} & T/C  \\   \hline
 31 & ERFPAT &Energy Flux, emitted, B.I.$\Rightarrow$Atoms&NATM
    & \si{\watt} & T/C  \\   \hline
 32 & EOTML &Energy Flux, incident, Molecules&NMOL
    & \si{\watt} & T/C  \\    \hline
 33 & ERFAML &Energy Flux, emitted, Ats.$\Rightarrow$Molecules&NMOL
    & \si{\watt} & T/C   \\   \hline
 34 & ERFMML &Energy Flux, emitted, Mls.$\Rightarrow$Molecules&NMOL
    & \si{\watt} & T/C   \\   \hline
 35 & ERFIML &Energy Flux, emitted, T.I.$\Rightarrow$Molecules&NMOL
    & \si{\watt} & T/C   \\   \hline
 36 & ERFPHML &Energy Flux, emitted, Pht.$\Rightarrow$Molecules&NMOL
    & \si{\watt} & T/C   \\   \hline
 37 & ERFPML &Energy Flux, emitted, B.I.$\Rightarrow$Molecules&NMOL
    & \si{\watt} & T/C   \\   \hline
 38 & EOTIO &Energy Flux, incident, Test Ions&NION
    & \si{\watt} & T/C   \\   \hline
 39 & ERFAIO & Energy Flux, emitted, Ats.$\Rightarrow$Test Ions & NION
    & \si{\watt} & T/C   \\   \hline
 40 & ERFMIO & Energy Flux, emitted, Mls.$\Rightarrow$Test Ions & NION
    & \si{\watt} & T/C   \\   \hline
 41 & ERFIIO & Energy Flux, emitted, T.I.$\Rightarrow$Test Ions & NION
    & \si{\watt} & T/C   \\   \hline
 42 & ERFPHIO & Energy Flux, emitted, Pht.$\Rightarrow$Test Ions & NION
    & \si{\watt} & T/C   \\   \hline
 43 & ERFPIO & Energy Flux, emitted, B.I.$\Rightarrow$Test Ions & NION
    & \si{\watt} & T/C   \\   \hline
 44 & EOTPHT &Energy Flux, incident, Photons&NPHOT
    & \si{\watt} & T/C   \\   \hline
 45 & ERFAPHT & Energy Flux, emitted, Ats.$\Rightarrow$Photons & NPHOT
    & \si{\watt} & T/C   \\   \hline
 46 & ERFMPHT & Energy Flux, emitted, Mls.$\Rightarrow$Photons & NPHOT
    & \si{\watt} & T/C   \\   \hline
 47 & ERFIPHT & Energy Flux, emitted, T.I.$\Rightarrow$Photons & NPHOT
    & \si{\watt} & T/C   \\   \hline
 48 & ERFPHPHT & Energy Flux, emitted, Pht.$\Rightarrow$Photons & NPHOT
    & \si{\watt} & T/C   \\   \hline
 49 & ERFPPHT & Energy Flux, emitted, B.I.$\Rightarrow$Photons & NPHOT
    & \si{\watt} & T/C   \\   \hline
 50 & EOTPL &Energy Flux, incident, Bulk Ions&NPLS
    & \si{\watt} & T/C   \\   \hline\hline
 51 & SPTAT &Sputtered Flux, by incident Atoms&NATM
    & \si{\ampere} & T/C  \\ \hline
 52 & SPTML &Sputtered Flux, by incident Molecules&NMOL
    & \si{\ampere} & T/C  \\  \hline
 53 & SPTIO &Sputtered Flux, by incident Test Ions&NION
    & \si{\ampere} & T/C  \\    \hline
 54 & SPTPHT &Sputtered Flux, by incident Photons&NPHOT
    & \si{\ampere} & T/C  \\    \hline
 55 & SPTPL &Sputtered Flux, by incident Bulk Ions&NPLS
    & \si{\ampere} & T/C  \\  \hline
 56 & SPTTOT &Sputtered Flux, total&1
    & \si{\ampere} & T/C  \\  \hline\hline
 57 & ADDS & Additional Surface Tally & NADS
    & Input & T/C              \\  \hline
 58 & ALGS & Algebraic expression in surface averaged tallies&NALS
    & Input &                       \\  \hline
 59 & SPUMP & Pumped flux, by species   &NSPZ
    & \si{\ampere} &                       \\  \hline
\end{longtable}
}

\textbf{Note:}\\
i) The tallies listed here are two dimensional arrays. The 2nd index I2
is the number of the surface or the surface segment.
 For
I2=1,...,NLIMI the tallies correspond to the ``Additional Surfaces",
(input block 3B).\\
ii)
For I2=NLIM+1,NLIM+NSTSI*NGITT the data correspond to the
``Non-default Standard Surfaces" (input block 3A).
On each such surface of block 3A,
there is a spatial resolution with up to NGITT
surface segments, depending upon the standard grid dimensionality.\\
iii) by default the sputter tallies are type and species resolved with respect to incident type and species.
E.g. the sputtered fluxes SPTAT(2) is the total sputtered flux (all species), sputtered by incident atoms of species IATM=2.
Other conventions regarding sputter tallies are used in EIRENE versions operated at ITER-IO (e.g. SOLPS4.x), see general description above.\\
iv) the sputter tallies have been revised at revision svn 360, Jan 2014, to allow full resolution with respect to both: incident type (for scaling)
and (emitted) sputtered particle type and species, see \cref{tab3a}.

%\end{table}

\newpage
\clearpage
{\footnotesize
\begin{longtable}{|r|l|l|l|l|c|}
\caption{\textbf{Sputter tallies, new version, svn 360 ff, 2014,\\
Output, Module: CESTIM}}
\label{tab3a}
\\
\hline

 \textbf{No} & \textbf{Name} & \textbf{Macroscopic quantity}
&
   \textbf{1.Dim.}
  & \textbf{Units} & \textbf{Estim.}   \\
\hline\hline
\endfirsthead
\caption[]{(continued)}\\
\hline
 \textbf{No} & \textbf{Name} & \textbf{Macroscopic quantity}
 & \textbf{1.Dim.}
 & \textbf{Units} & \textbf{Estim.} \\
\hline\hline
\endhead
 51 & SPTAAT &Sputtered Flux, Atoms, by incident Atoms&NATM
    & \si{\ampere} & T/C  \\ \hline
 52 & SPTAML &Sputtered Flux, Molecules, by incident Atoms&NMOL
    & \si{\ampere} & T/C  \\  \hline
 53 & SPTAIO &Sputtered Flux, Test Ions, by incident Atoms&NION
    & \si{\ampere} & T/C  \\    \hline
 54 & SPTAPHT &Sputtered Flux, Photons, by incident Atoms&NPHOT
    & \si{\ampere} & T/C  \\    \hline
 55 & SPTAPL &Sputtered Flux, Bulk Ions, by incident Atoms&NPLS
    & \si{\ampere} & T/C  \\  \hline\hline
 56 & SPTMAT &Sputtered Flux, Atoms, by incident Molecules&NATM
    & \si{\ampere} & T/C  \\ \hline
 57 & SPTMML &Sputtered Flux, Molecules, by incident Molecules&NMOL
    & \si{\ampere} & T/C  \\  \hline
 58 & SPTMIO &Sputtered Flux, Test Ions, by incident Molecules&NION
    & \si{\ampere} & T/C  \\    \hline
 59 & SPTMPHT &Sputtered Flux, Photons, by incident Molecules&NPHOT
    & \si{\ampere} & T/C  \\    \hline
 60 & SPTMPL &Sputtered Flux, Bulk Ions, by incident Molecules&NPLS
    & \si{\ampere} & T/C  \\  \hline\hline
 61 & SPTIAT &Sputtered Flux, Atoms, by incident Test Ions&NATM
    & \si{\ampere} & T/C  \\ \hline
 62 & SPTIML &Sputtered Flux, Molecules, by incident Test Ions&NMOL
    & \si{\ampere} & T/C  \\  \hline
 63 & SPTIIO &Sputtered Flux, Test Ions, by incident Test Ions&NION
    & \si{\ampere} & T/C  \\    \hline
 64 & SPTIPHT &Sputtered Flux, Photons, by incident Test Ions&NPHOT
    & \si{\ampere} & T/C  \\    \hline
 65 & SPTIPL &Sputtered Flux, Bulk Ions, by incident Test Ions&NPLS
    & \si{\ampere} & T/C  \\  \hline\hline
 66 & SPTPHAT &Sputtered Flux, Atoms, by incident Photons&NATM
    & \si{\ampere} & T/C  \\ \hline
 67 & SPTPHML &Sputtered Flux, Molecules, by incident Photons&NMOL
    & \si{\ampere} & T/C  \\  \hline
 68 & SPTPHIO &Sputtered Flux, Test Ions, by incident Photons&NION
    & \si{\ampere} & T/C  \\    \hline
 69 & SPTPHPHT &Sputtered Flux, Photons, by incident Photons&NPHOT
    & \si{\ampere} & T/C  \\    \hline
 70 & SPTPHPL &Sputtered Flux, Bulk Ions, by incident Photons&NPLS
    & \si{\ampere} & T/C  \\  \hline \hline
 71 & SPTPAT &Sputtered Flux, Atoms,  by incident Bulk Ions&NATM
    & \si{\ampere} & T/C  \\ \hline
 72 & SPTPML &Sputtered Flux, Molecules, by incident Bulk Ions&NMOL
    & \si{\ampere} & T/C  \\  \hline
 73 & SPTPIO &Sputtered Flux, Test Ions, by incident Bulk Ions&NION
    & \si{\ampere} & T/C  \\    \hline
 74 & SPTPPHT &Sputtered Flux, Photons, by incident Bulk Ions&NPHOT
    & \si{\ampere} & T/C  \\    \hline
 75 & SPTPPL &Sputtered Flux, Bulk Ions, by incident Bulk Ions&NPLS
    & \si{\ampere} & T/C  \\  \hline \hline
 76 & SPTATOP &Sputtered Flux, by incident Atoms & 1
    & \si{\ampere} & T/C  \\ \hline
 77 & SPTMTOT &Sputtered Flux, by incident Molecules& 1
    & \si{\ampere} & T/C  \\  \hline
 78 & SPTITOP &Sputtered Flux, by incident Test Ions& 1
    & \si{\ampere} & T/C  \\    \hline
 79 & SPTPHTOT &Sputtered Flux, by incident Photons& 1
    & \si{\ampere} & T/C  \\    \hline
 80 & SPTPLTOT &Sputtered Flux, by incident Bulk Ions& 1
    & \si{\ampere} & T/C  \\  \hline \hline
 81 & SPTTOT &Sputtered Flux, total&1
    & \si{\ampere} & T/C  \\  \hline\hline
 82 & ADDS & Additional Surface Tally & NADS
    & Input & T/C              \\  \hline
 83 & ALGS & Algebraic expression in surface averaged tallies&NALS
    & Input &                       \\  \hline
 84 & SPUMP & Pumped flux, by species   &NSPZ
    & \si{\ampere} &                       \\  \hline
\end{longtable}
}
%\end{table}

\textbf{Note: }\\
This extended set of sputter tallies has now species naming conventions SPT[A][BC] fully identical to those for volumetric source rates,
as well as those for emitted particle and energy fluxes from surfaces.
The full terminological equivalence leads to presence of quite strange sputter tallies (e.g. bulk ion fluxes sputtered by incident photons,
tally 70), but redundant tallies (and related storage) are automatically removed from a run.
Physically irrelevant tallies are listed here only to maintain complete symmetry in notation.
\clearpage
\newpage
\subsection{old version, w/o photon gas tallies (Eirene\texorpdfstring{$_{2001}$}{2001} and
older)}\label{sec.tables-old}
In these older versions photon test particle tallies are not available, nor are the
momentum source rates MAPL,.... and test particle momentum fluxes VXDEN...,VYDEN,...VZDEN.
In these versions such tallies may still have been available, but then in problem specific tallies
COPV (interfacing to particular plasma fluid codes) or in tally BGKV (internal iterations due to neutral-neutral collisions
in \ac{BGK} approximation).

\medskip
With respect to surface averaged tallies (in and outgoing surface fluxes, sputter fluxes,
etc...) there have been major revisions in versions
2014 or later, see ``sputter tallies, new version, svn 360 ff, 2014" in
\cref{sec.tables-new}.

A particular modification (A. Kukushkin) of the sputter tallies has been used
in many B2-EIRENE applications,
in versions SOLPS4.x, see under \cref{sec2.14} in which problem specific issues regarding EIRENE
interfacing to the particular external plasma codes B2, B2.5, etc. are described.

\begin{table}[ht]
\caption{\textbf{Input Tallies for Background, Input, Common: COMUSR}}
\bigskip
\label{tab1old} {\footnotesize
\begin{tabular} {|r|l|l|l|l|} \hline

 \textbf{No} & \textbf{Name} & \textbf{Macroscopic quantity}
 &  \textbf{1.Dim.}
  & \textbf{Units}   \\ \hline

 -1 & TEIN & Plasma Temperature, Electrons & 1 & \si{\electronvolt}   \\ \hline
 -2 & TIIN & Plasma Temperature, Bulk Ions & NPLS & \si{\electronvolt}   \\ \hline
 -3 & DEIN$^*$ & Plasma Density, Electrons & 1 & \si{\per\cubic\centi\metre}   \\ \hline
 -4 & DIIN & Plasma Density, Bulk Ions & NPLS & \si{\per\cubic\centi\metre}  \\ \hline
 -5 & VXIN & Plasma Drift Velocity, x-component, Bulk Ions&NPLS  & \si{\centi\metre\per\second}  \\ \hline
 -6 & VYIN & Plasma Drift Velocity, y-component, Bulk Ions&NPLS
     & \si{\centi\metre\per\second}                                            \\ \hline
 -7 & VZIN & Plasma Drift Velocity, z-component, Bulk Ions&NPLS
     & \si{\centi\metre\per\second}                                            \\ \hline
 -8 & BXIN & Magnetic field unit vector, x-component & 1 & /
    \\ \hline
 -9 & BYIN & Magnetic field unit vector, y-component & 1 & /
    \\ \hline
 -10 & BZIN & Magnetic field unit vector, z-component & 1 & /
    \\ \hline
 -11 & BFIN & Magnetic field strength & 1 &  Tesla
    \\ \hline
 -12 & ADIN & Additional input tally & NAIN & see INFCOP
    \\ \hline
 -13 & EDRIFT$^*$ & Kinetic energy in drift motion, Bulk Ions &NPLS & \si{\electronvolt}
    \\ \hline
 -14 & VOL & Zone Volume & 1 & \si{\cubic\centi\metre}
    \\ \hline
 -15 & WGHT$^{**}$ & Space and species dependent importance & NSPCMC & $ 1 $
    \\ \hline
 -22   &      POT   &      electric potential  & 1  &         V  \\ \hline
\end{tabular}
}

\textbf{Note:}
We distinguish between so called primary input tallies, and derived input tallies.
The latter are automatically computed from the former by EIRENE.


Primary input tallies are set or transferred from external data sources via subr. PLASMA.f,
derived input tallies are set via subroutine PLASMA$\_$DERIV.f.
\begin{itemize}
\item DEIN and EDRIFT are ``derived" input tallies,
internally computed from the ion densities (quasi-neutrality), and
the flow field, respectively.
\item  WGHT space species  weight window arrays.  Currently not in use
\item  Magnetic field input is needed only for photon transport with Zeeman splitted line shapes, or
in case of charged particle transport (Subr. FOLION). Only quite rudimentary preprogrammed input options
for magnetic field are available, usually it is expected to receive B-field information from external files,
see INDPRO(5) options for magnetic field tallies in input block 5 (\cref{sec2.5}).

\end{itemize}
\end{table}





\newpage
\begin{table}[t]
\caption{\textbf{Volume Averaged Tallies, Output, Common: CESTIM}}
\label{tab2old} {\footnotesize
\begin{tabular}{|r|l|l|l|l|c|} \hline
 \textbf{No} & \textbf{Name} & \textbf{Macroscopic quantity}
 & \textbf{1.Dim.}
 & \textbf{Units} & \textbf{Estim.} \\ \hline
 1 & PDENA & Particle density, Atoms & NATM & \si{\per\cubic\centi\metre} & T1
 \\ \hline
 2 & PDENM & Particle density, Molecules & NMOL & \si{\per\cubic\centi\metre} & T2
 \\ \hline
 3 & PDENI & Particle density Test Ions & NION  & \si{\per\cubic\centi\metre} & T3
 \\ \hline
 4 & EDENA & Energy density, Atoms & NATM       & \si{\electronvolt}*\si{\per\cubic\centi\metre} & T4
 \\ \hline
 5 & EDENM & Energy density, Molecules & NMOL   & \si{\electronvolt}*\si{\per\cubic\centi\metre} & T5
 \\ \hline
 6 & EDENI & Energy density, Test Ions & NION   & \si{\electronvolt}*\si{\per\cubic\centi\metre} & T6
 \\ \hline
 7 & PAEL & Particle Source (Electrons) from atom-plasma coll.
   & 1 & \si{\ampere}* \si{\per\cubic\centi\metre} & T7  \\ \hline
 8 & PAAT & Particle Source (Atoms) from atom-plasma coll.
   & NATM & \si{\ampere}*\si{\per\cubic\centi\metre} & T8 \\ \hline
 9 & PAML & Particle Source (Molecules) from atom-plasma coll.
   & NMOL & \si{\ampere}*\si{\per\cubic\centi\metre} & T9 \\  \hline
 10 & PAIO & Particle Source (Test Ions) from atom-plasma coll.
    & NION & \si{\ampere}*\si{\per\cubic\centi\metre} & T10  \\  \hline
 11 & PAPL & Particle Source (Bulk Ions) from atom-plasma coll.
    & NPLS & \si{\ampere}*\si{\per\cubic\centi\metre} & T11  \\  \hline
 12 & PMEL &Particle Source (Electrons) from molecule-plasma coll.
    & 1 & \si{\ampere}*\si{\per\cubic\centi\metre} & T12     \\  \hline
 13 & PMAT &Particle Source (Atoms) from molecule-plasma coll.
    & NATM &\si{\ampere}*\si{\per\cubic\centi\metre} & T13  \\  \hline
 14 & PMML &Particle Source (Molecules) from molecule-plasma coll.
    & NMOL &\si{\ampere}*\si{\per\cubic\centi\metre} & T14   \\  \hline
 15 & PMIO &Particle Source (Test Ions) from molecule-plasma coll.
    &NION&\si{\ampere}*\si{\cubic\centi\metre} & T15     \\   \hline
 16 & PMPL &Particle Source (Bulk Ions) from molecule-plasma coll.
    &NPLS&\si{\ampere}*\si{\per\cubic\centi\metre} & T16    \\    \hline
 17 & PIEL &Particle Source (Electrons) from test ion-plasma coll.
    &1&\si{\ampere}*\si{\per\cubic\centi\metre} & T17       \\   \hline
 18 & PIAT &Particle Source (Atoms) from test ion-plasma coll.
    & NATM&\si{\ampere}*\si{\per\cubic\centi\metre} & T18    \\  \hline
 19 & PIML &Particle Source (Molecules) from test ion-plasma coll.
    & NMOL&\si{\ampere}*\si{\per\cubic\centi\metre} & T19    \\  \hline
 20 & PIIO &Particle Source (Test Ions) from test ion-plasma coll.
    & NION&\si{\ampere}*\si{\per\cubic\centi\metre} & T20    \\   \hline
 21 & PIPL &Particle Source (Bulk Ions) from test ion-plasma coll.
    & NPLS&\si{\ampere}*\si{\per\cubic\centi\metre} & T21    \\   \hline
 22 & EAEL&Energy Source (Electrons) from atom-plasma coll.
    & 1 & \si{\watt}*\si{\per\cubic\centi\metre} & T22      \\  \hline
 23 & EAAT&Energy Source (Atoms) from atom-plasma coll.
    & 1 & \si{\watt}*\si{\per\cubic\centi\metre} & T23      \\   \hline
 24 & EAML&Energy Source (Molecules) from atom-plasma coll.
    & 1 & \si{\watt}*\si{\per\cubic\centi\metre} & T24      \\    \hline
 25 & EAIO&Energy Source (Test Ions) from atom-plasma coll.
    & 1 & \si{\watt}*\si{\per\cubic\centi\metre} & T25      \\ \hline
 26 & EAPL&Energy Source (Bulk Ions) from atom-plasma coll.
    & 1 & \si{\watt}*\si{\per\cubic\centi\metre} & T26   \\  \hline
 27 & EMEL&Energy Source (Electrons) from molecule-plasma coll.
    & 1 & \si{\watt}*\si{\per\cubic\centi\metre} & T27   \\   \hline
 28 & EMAT&Energy Source (Atoms) from molecule-plasma coll.
    & 1 & \si{\watt}*\si{\per\cubic\centi\metre} & T28  \\   \hline
 29 & EMML&Energy Source (Molecules) from molecule-plasma coll.
    & 1 & \si{\watt}*\si{\per\cubic\centi\metre} & T29  \\   \hline
 30 & EMIO&Energy Source (Test Ions) from molecule-plasma coll.
    & 1 & \si{\watt}*\si{\per\cubic\centi\metre} & T30  \\    \hline
 31 & EMPL & Energy Source (Bulk Ions) from molecule-plasma coll.
    & 1& \si{\watt}*\si{\per\cubic\centi\metre} & T31  \\  \hline
 32 & EIEL & Energy Source (Electrons) from test ion-plasma coll.
    & 1& \si{\watt}*\si{\per\cubic\centi\metre} & T32   \\  \hline
 33 & EIAT & Energy Source (Atoms) from test ion-plasma coll.
    & 1& \si{\watt}*\si{\per\cubic\centi\metre} & T33   \\  \hline
 34 & EIML & Energy Source (Molecules) from test ion-plasma coll.
    & 1& \si{\watt}*\si{\per\cubic\centi\metre} & T34   \\  \hline
 35 & EIIO & Energy Source (Test Ions) from test ion-plasma coll.
    & 1& \si{\watt}*\si{\per\cubic\centi\metre} & T35   \\  \hline
 36 & EIPL & Energy Source (Bulk Ions) from test ion-plasma coll.
    & 1& \si{\watt}*\si{\per\cubic\centi\metre} & T36   \\  \hline
 37 & ADDV & Additional volume av. Tally, Track-length estimated
    & NADV & see UPTUSR & TRL      \\  \hline
 38 & COLV & Additional volume av. Tally, Collision estimated
    & NCLV & see UPCUSR & COL        \\  \hline
 39 & SNPV & Additional volume av. Tally, Snapshot estimated
    & NSNV & see UPSUSR & SNP         \\  \hline
 40 & COPV & Tallies for coupling to ext. code, see: Subr. UPTCOP
    & NCPV & see UPTCOP  & COP  \\  \hline
 41 & BGKV & Volume averaged tallies for \acs{BGK}-self-collision terms
    & NBGV & see BGK & T41        \\  \hline
 42 & ALGV & Algebraic expression in volume averaged tallies
    & NALV & Input & ALG        \\  \hline
 43 & PGENA & Generation limit, Atoms & NATM & \si{\ampere}*\si{\per\cubic\centi\metre} & T1
 \\ \hline
 44 & PGENM & Generation limit, Molecules & NMOL & \si{\ampere}*\si{\per\cubic\centi\metre} & T1
 \\ \hline
 45 & PGENI & Generation limit, Test ions & NION & \si{\ampere}*\si{\per\cubic\centi\metre} & T1
 \\ \hline
 46 & EGENA & dito, Energy flux, Atoms & NATM & \si{\watt}*\si{\per\cubic\centi\metre} & T1
 \\ \hline
 47 & EGENM & dito, Energy flux, Molecules & NMOL
& \si{\watt}*\si{\per\cubic\centi\metre} & T1
 \\ \hline
 48 & EGENI & dito, Energy flux, Test ions & NION & \si{\watt}*\si{\per\cubic\centi\metre} & T1
 \\ \hline
 49 & VGENA & dito, momentum flux, Atoms & NATM & as COP
& T1
 \\ \hline
 50 & VGENM & dito, momentum flux, Molecules & NMOL & as
COP & T1
 \\ \hline
 51 & VGENI & dito, Momentum flux, Test ions &
NION & as COP & T1
 \\ \hline
\end{tabular}
}
\textbf{Note:} $\underline{Source} $ means: if the sign is positive,
it is a gain for the specified type of particles; if it is
negative, it is a loss.
\end{table}

\newpage
\begin{table}[t]
\caption{\textbf{Surface Averaged Tallies, Output, Common: CESTIM}}
\bigskip
\label{tab3old} {\footnotesize
\begin{tabular}{|r|l|l|l|l|c|} \hline

 \textbf{No} & \textbf{Name} & \textbf{Macroscopic quantity}
&
   \textbf{1.Dim.}
  & \textbf{Units} & \textbf{Estim.}   \\ \hline
 1 & POTAT&Particle Flux, incident, Atoms&NATM
    & \si{\ampere} & T/C  \\ \hline
 2 &PRFAAT&Particle Flux, emitted, Ats.$\Rightarrow$Atoms&NATM
    & \si{\ampere} & T/C  \\  \hline
 3 &PRFMAT&Particle Flux, emitted, Mls.$\Rightarrow$Atoms&NATM
    & \si{\ampere} & T/C  \\  \hline
 4 &PRFIAT&Particle Flux, emitted, T.I.$\Rightarrow$Atoms&NATM
    & \si{\ampere} & T/C  \\  \hline
 5 &PRFPAT&Particle Flux, emitted, B.I.$\Rightarrow$Atoms&NATM
    & \si{\ampere} & T/C  \\  \hline
 6 &POTML&Particle Flux, incident, Molecules&NMOL
    & \si{\ampere} & T/C  \\  \hline
 7 &PRFAML&Particle Flux, emitted, Ats.$\Rightarrow$Molecules&NMOL
    & \si{\ampere} & T/C  \\   \hline
 8 &PRFMML&Particle Flux, emitted, Mls.$\Rightarrow$Molecules&NMOL
    & \si{\ampere} & T/C  \\   \hline
 9 &PRFIML&Particle Flux, emitted, T.I.$\Rightarrow$Molecules&NMOL
    & \si{\ampere} & T/C  \\   \hline
 10 &PRFPML&Particle Flux, emitted, B.I.$\Rightarrow$Molecules&NMOL
    & \si{\ampere} & T/C  \\   \hline
 11 &POTIO&Particle Flux, incident, Test Ions&NION
    & \si{\ampere} & T/C  \\    \hline
 12 &PRFAIO&Particle Flux, emitted, Ats.$\Rightarrow$Test Ions&NION
    & \si{\ampere} & T/C  \\    \hline
 13 &PRFMIO&Particle Flux, emitted, Mls.$\Rightarrow$Test Ions&NION
    & \si{\ampere} & T/C  \\    \hline
 14 &PRFIIO&Particle Flux, emitted, T.I.$\Rightarrow$Test Ions&NION
    & \si{\ampere} & T/C  \\    \hline
 15 &PRFPIO&Particle Flux, emitted, B.I.$\Rightarrow$Test Ions&NION
    & \si{\ampere} & T/C  \\    \hline
 16 &POTPL&Particle Flux, incident, Bulk Ions&NPLS
    & \si{\ampere} & T/C  \\    \hline
 17 &EOTAT&Energy Flux, incident, Atoms&NATM
    & \si{\watt} & T/C  \\   \hline
 18 &ERFAAT&Energy Flux, emitted, Ats.$\Rightarrow$Atoms&NATM
    & \si{\watt} & T/C  \\   \hline
 19 &ERFMAT&Energy Flux, emitted, Mls.$\Rightarrow$Atoms&NATM
    & \si{\watt} & T/C  \\   \hline
 20 &ERFIAT&Energy Flux, emitted, T.I.$\Rightarrow$Atoms&NATM
    & \si{\watt} & T/C  \\   \hline
 21 &ERFPAT&Energy Flux, emitted, B.I.$\Rightarrow$Atoms&NATM
    & \si{\watt} & T/C  \\   \hline
 22 &EOTML&Energy Flux, incident, Molecules&NMOL
    & \si{\watt} & T/C  \\    \hline
 23 &ERFAML&Energy Flux, emitted, Ats.$\Rightarrow$Molecules&NMOL
    & \si{\watt} & T/C   \\   \hline
 24 &ERFMML&Energy Flux, emitted, Mls.$\Rightarrow$Molecules&NMOL
    & \si{\watt} & T/C   \\   \hline
 25 &ERFIML&Energy Flux, emitted, T.I.$\Rightarrow$Molecules&NMOL
    & \si{\watt} & T/C   \\   \hline
 26 &ERFPML&Energy Flux, emitted, B.I.$\Rightarrow$Molecules&NMOL
    & \si{\watt} & T/C   \\   \hline
 27 &EOTIO&Energy Flux, incident, Test Ions&NION
    & \si{\watt} & T/C   \\   \hline
 28 & ERFAIO & Energy Flux, emitted, Ats.$\Rightarrow$Test Ions & NION
    & \si{\watt} & T/C   \\   \hline
 29 & ERFMIO & Energy Flux, emitted, Mls.$\Rightarrow$Test Ions & NION
    & \si{\watt} & T/C   \\   \hline
 30 & ERFIIO & Energy Flux, emitted, T.I.$\Rightarrow$Test Ions & NION
    & \si{\watt} & T/C   \\   \hline
 31 & ERFPIO & Energy Flux, emitted, B.I.$\Rightarrow$Test Ions & NION
    & \si{\watt} & T/C   \\   \hline
 32 &EOTPL&Energy Flux, incident, Bulk Ions&NPLS
    & \si{\watt} & T/C   \\   \hline
 33 & SPTAT&Sputtered Flux, by incident Atoms&NATM
    & \si{\ampere} & T/C  \\ \hline
 34 & SPTML&Sputtered Flux, by incident Molecules&NMOL
    & \si{\ampere} & T/C  \\  \hline
 35 & SPTIO&Sputtered Flux, by incident Test Ions&NION
    & \si{\ampere} & T/C  \\    \hline
 36 & SPTPL&Sputtered Flux, by incident Bulk Ions&NPLS
    & \si{\ampere} & T/C  \\  \hline
 37 & SPTTOT&Sputtered Flux, total&1
    & \si{\ampere} & T/C  \\  \hline
 38 & ADDS & Additional Surface Tally & NADS
    & Input & T/C              \\  \hline
 39 & ALGS & Algebraic expression in surface averaged tallies&NALS
    & Input &                       \\  \hline
 40 & SPUMP & Pumped flux, by species   & NSPZ
    & \si{\ampere} &                       \\  \hline
\end{tabular}
}
\\
\textbf{Note:} The tallies listed here are two dimensional arrays.
The 2nd index I2 is the number of the surface or the surface
segment.
 For
I2=1,...,NLIMI the tallies correspond to the ``Additional
Surfaces", (input block 3B).

For I2=NLIM+1,NLIM+NSTSI*NGITT the data correspond to the
``Non-default Standard Surfaces" (input block 3A). On each such
surface, there is a spatial resolution with up to NGITT surface
segments.

\end{table}
